\begin{resumo}[Abstract]
\begin{otherlanguage*}{english}
Monitoring the training of deep neural networks currently relies on a labelled
validation set that is representative and sufficiently large --- a condition
often unattainable in medical imaging, where annotation is costly and requires
expert knowledge. This dissertation describes how statistical complexity
measures computed exclusively from network weights, without access to any
label, behave throughout training. Two measures are considered: the
Lopez-Ruiz--Mancini--Calbet (LMC) complexity, defined as the product of Shannon
entropy and the statistical distance from equilibrium, and sample entropy
(SampEn), computed at every epoch over the dense-layer weights of
ResNet-18 networks trained on two medical image datasets, MedNIST and
PathMNIST. The trajectories are compared across control runs, runs with a
subsampled training set --- with and without matching the number of
optimisation steps to the control --- and runs with white and pink noise added
to the training images. In both datasets, LMC rises and SampEn falls during
training. With few optimisation steps, subsampling reduces the final LMC to
between 0.13 and 0.54 times that of the control. On MedNIST this difference
practically vanishes when the steps are matched; on PathMNIST, where the
network memorises the reduced subset, it persists in every statistic. Image
noise, although it degrades performance --- on PathMNIST, test accuracy drops
from 0.88 to 0.48 ---, does not change the measures on either dataset. The
dense-layer measures therefore follow the number of optimisation steps and,
when the network memorises, the amount of data, but do not register image
degradation.

\vspace{\onelineskip}
\noindent
\textbf{Keywords}: Statistical complexity; Sample entropy; Deep learning;
Medical imaging; LMC complexity.
\end{otherlanguage*}
\end{resumo}
